% Options for packages loaded elsewhere
\PassOptionsToPackage{unicode}{hyperref}
\PassOptionsToPackage{hyphens}{url}
\PassOptionsToPackage{dvipsnames,svgnames,x11names}{xcolor}
\documentclass[
  10pt,
  fontset=fandol]{ctexart}
\usepackage{xcolor}
\usepackage[margin=16mm]{geometry}
\usepackage{amsmath,amssymb}
\setcounter{secnumdepth}{-\maxdimen} % remove section numbering
\usepackage{iftex}
\ifPDFTeX
  \usepackage[T1]{fontenc}
  \usepackage[utf8]{inputenc}
  \usepackage{textcomp} % provide euro and other symbols
\else % if luatex or xetex
  \usepackage{unicode-math} % this also loads fontspec
  \defaultfontfeatures{Scale=MatchLowercase}
  \defaultfontfeatures[\rmfamily]{Ligatures=TeX,Scale=1}
\fi
\usepackage{lmodern}
\ifPDFTeX\else
  % xetex/luatex font selection
\fi
% Use upquote if available, for straight quotes in verbatim environments
\IfFileExists{upquote.sty}{\usepackage{upquote}}{}
\IfFileExists{microtype.sty}{% use microtype if available
  \usepackage[]{microtype}
  \UseMicrotypeSet[protrusion]{basicmath} % disable protrusion for tt fonts
}{}
\makeatletter
\@ifundefined{KOMAClassName}{% if non-KOMA class
  \IfFileExists{parskip.sty}{%
    \usepackage{parskip}
  }{% else
    \setlength{\parindent}{0pt}
    \setlength{\parskip}{6pt plus 2pt minus 1pt}}
}{% if KOMA class
  \KOMAoptions{parskip=half}}
\makeatother
\setlength{\emergencystretch}{3em} % prevent overfull lines
\providecommand{\tightlist}{%
  \setlength{\itemsep}{0pt}\setlength{\parskip}{0pt}}
\usepackage{amsmath,amssymb,mathtools,needspace}
\xeCJKDeclareCharClass{CJK}{"2032,"0400->"04FF,"2160->"216B,"2460->"2473,"25A0,"25C7,"25CB,"25CF}
\IfFontExistsTF{Arial Unicode MS}{\xeCJKsetup{AutoFallBack=true}\setCJKfallbackfamilyfont{\CJKrmdefault}{Arial Unicode MS}}{}
\setcounter{secnumdepth}{0}
\setlength{\emergencystretch}{2em}
\allowdisplaybreaks
\usepackage{bookmark}
\IfFileExists{xurl.sty}{\usepackage{xurl}}{} % add URL line breaks if available
\urlstyle{same}
\hypersetup{
  pdftitle={Walsh 阵的演化生成},
  colorlinks=true,
  linkcolor={Maroon},
  filecolor={Maroon},
  citecolor={Blue},
  urlcolor={Blue},
  pdfcreator={LaTeX via pandoc}}

\title{Walsh 阵的演化生成}
\author{}
\date{}

\begin{document}
\maketitle

\subsubsection{10.3.2　Walsh 阵的演化生成}\label{walsh-ux9635ux7684ux6f14ux5316ux751fux6210}

首先考察表 10.1 中镜像行复制的演化方式。考察某个方阵 \(A\)，用 \(A(i)\) 表示其第 \(i\) 行，对 \(A(i)\) 施行偶复制与奇复制，分别生成向量 \([A(i)\mathrel{\vdots}\ddot A(i)]\) 与 \([A(i)\mathrel{\vdots}\dot A(i)]\)。

例如，设 \(A(i)=[+\ -]\)，则

\[
[A(i)\mathrel{\vdots}\ddot A(i)]=[+\ -\mathrel{\vdots}-\ +],
\]

\[
[A(i)\mathrel{\vdots}\dot A(i)]=[+\ -\mathrel{\vdots}+\ -].
\]

进一步，若 \(A(i)=[+\ -\mathrel{\vdots}+\ -]\)，则

\[
[A(i)\mathrel{\vdots}\ddot A(i)]=[+\ -\ +\ -\mathrel{\vdots}-\ +\ -\ +],
\]

\[
[A(i)\mathrel{\vdots}\dot A(i)]=[+\ -\ +\ -\mathrel{\vdots}+\ -\ +\ -].
\]

如果对方阵 \(A\) 的每一行先后施行偶复制与奇复制两种复制手续，即可生成一个阶数倍增的方阵 \(B\)，这种演化手续称作\textbf{镜像行复制}，即

\href{../../source-images/pdf-267.jpeg}{原页}

\[
A=\begin{pmatrix}\vdots\\A(i)\\\vdots\end{pmatrix}
\longrightarrow B=\begin{pmatrix}
\vdots\\
A(i)\mathrel{\vdots}\ddot A(i)\\
A(i)\mathrel{\vdots}\dot A(i)\\
\vdots
\end{pmatrix}.
\]

人们自然会问，如果对方波 \([+]\) 反复施行镜像行复制的演化手续，使其阶数逐步倍增，将会生成什么样的方阵序列呢？

1 阶方阵 \([+]\) 仅有一行（一列），对它施行偶复制与奇复制，分别生成 \([+\mathrel{\vdots}+]\) 与 \([+\mathrel{\vdots}-]\)，两者合成在一起，结果生成一个 2 阶方阵

\[
[+]\longrightarrow\begin{pmatrix}+&+\\+&-\end{pmatrix}.
\]

对所生成的 2 阶方阵的两行 \([+\ +]\) 与 \([+\ -]\) 分别施行镜像复制的偶复制与奇复制，进一步生成一个 4 阶方阵

\[
\begin{pmatrix}+&+\\+&-\end{pmatrix}\longrightarrow
\begin{pmatrix}
+&+&+&+\\
+&+&-&-\\
+&-&-&+\\
+&-&+&-
\end{pmatrix}.
\]

继续对所生成的 4 阶方阵施行镜像行复制，获得如下 8 阶方阵，即

\[
\begin{pmatrix}
+&+&+&+\\
+&+&-&-\\
+&-&-&+\\
+&-&+&-
\end{pmatrix}\longrightarrow
\begin{pmatrix}
+&+&+&+&+&+&+&+\\
+&+&+&+&-&-&-&-\\
+&+&-&-&-&-&+&+\\
+&+&-&-&+&+&-&-\\
+&-&-&+&+&-&-&+\\
+&-&-&+&-&+&+&-\\
+&-&+&-&-&+&-&+\\
+&-&+&-&+&-&+&-
\end{pmatrix}.
\]

上述方阵与 10.2.2 小节列出的 Walsh 方阵相比较，两者完全一致。可以证明，从方波 \([+]\) 出发，运用镜像行复制的演化技术可以生成 Walsh 方阵序列。这个 Walsh 方阵等价于 10.1.2 小节的原始定义。

Walsh 方阵有多种排序方式。镜像行复制生成的 Walsh 方阵特别称作 \textbf{Walsh 阵}。

\end{document}
